\chapter{Programm-Listings}\label{ch:robotik-programme}

% Die Originaldateien einlesen, damit Skript und Programme übereinstimmen.
% Die Listing-Einstellungen gelten nur innerhalb dieses Anhangs.
\begingroup
\lstset{
    language=C++,
    columns=fullflexible,
    keepspaces=true,
    breaklines=true,
    showspaces=false,
    showstringspaces=false,
    showtabs=false,
    tabsize=2,
    numbers=none,
    basicstyle=\ttfamily\footnotesize
}

\section{Distanzmessung mit Ultraschall}

\lstinputlisting[
    caption={\texttt{Intro\_Ultraschall.ino}},
    label={lst:ultraschall-intro}
]{FF/Robotik/Skript/Code/Intro_Ultraschall/Intro_Ultraschall.ino}

\lstinputlisting[
    caption={\texttt{ultraSonicTest1}},
    label={lst:ultraschall-roboter}
]{FF/Robotik/Skript/Code/ultraSonicTest_1/ultraSonicTest_1.ino}

\lstinputlisting[
    caption={\texttt{hinundher\_template.ino}},
    label={lst:hinundher-template}
]{FF/Robotik/Skript/Code/hinundher_template/hinundher_template.ino}

\section{Darstellung der Messwerte in Polarkoordinaten}

\lstinputlisting[
    language=Python,
    caption={\texttt{polarkoord-sketch}},
    label={lst:polarkoord-sketch}
]{FF/Robotik/Skript/Code/polarkoord_sketch.py}

\section{Distanzmessung mit Infrarot}

\lstinputlisting[
    caption={\texttt{infraRedTest}},
    label={lst:infrarot-test}
]{FF/Robotik/Skript/Code/infraRedTest/infraRedTest.ino}

\section{Berührungssensoren}

\lstinputlisting[
    caption={\texttt{TouchSensorTest.ino}},
    label={lst:touchsensor-test}
]{FF/Robotik/Skript/Code/TouchSensorTest/TouchSensorTest.ino}

\section{Lichtsensoren}

\lstinputlisting[
    caption={\texttt{LightSensorTest.ino}},
    label={lst:lichtsensor-test}
]{FF/Robotik/Skript/Code/LightSensorTest/LightSensorTest.ino}

\endgroup
